\subsection{Graph Partitioning Problem}

Given an undirected graph $G = (V, E)$ with vertex weights $w: V \to \mathbb{R}^+$, partition vertices into $k$ disjoint sets $\{P_1, P_2, \ldots, P_k\}$ that:

\begin{enumerate}
\item \textbf{Cover all vertices}: $\bigcup_{i=1}^k P_i = V$ and $P_i \cap P_j = \emptyset$ for $i \neq j$
\item \textbf{Balance weights}: $|w(P_i) - w_{target}| \leq \epsilon \cdot w_{target}$ where $w(P_i) = \sum_{v \in P_i} w(v)$ and $w_{target} = \frac{1}{k}\sum_{v \in V} w(v)$
\item \textbf{Minimize edge cut}: $\text{EdgeCut} = |\{(u,v) \in E : u \in P_i, v \in P_j, i \neq j\}|$
\end{enumerate}

This problem is NP-hard \citep{garey1979}. Practical algorithms use multilevel heuristics that coarsen the graph, partition the coarse graph, and refine the partition \citep{karypis1998}.

\subsection{METIS Partitioning Library}

METIS \citep{karypis1998metis} is the standard implementation for graph partitioning. It operates in three phases:

\textbf{Coarsening}: Iteratively merge vertices to create smaller graphs $G_0, G_1, \ldots, G_m$ where $G_0 = G$ and $|V_m| \ll |V_0|$.

\textbf{Initial Partitioning}: Partition coarsest graph $G_m$ into $k$ parts using greedy or spectral methods.

\textbf{Refinement}: Project partition back through $G_{m-1}, G_{m-2}, \ldots, G_0$, refining at each level using Kernighan-Lin (KL) or greedy algorithms to improve edge cut while maintaining balance.

METIS provides two partitioning strategies:

\subsection{Recursive Bisection (PartGraphRecursive)}

Algorithm: Partition $G$ into 2 parts. Recursively partition each part until $k$ partitions exist.

\begin{algorithm}[H]
\caption{Recursive Bisection}
\begin{algorithmic}[1]
\STATE \textbf{Input:} Graph $G$, target partitions $k$
\STATE \textbf{Output:} Partition assignment $P: V \to \{1, \ldots, k\}$
\STATE
\IF{$k = 1$}
    \RETURN single partition
\ENDIF
\STATE
\STATE Choose split: $k = k_L + k_R$ where $k_L, k_R \geq 1$
\STATE Set target weights: $w_L = \frac{k_L}{k}W$, $w_R = \frac{k_R}{k}W$ where $W = \sum_{v \in V} w(v)$
\STATE Partition $G$ into $(G_L, G_R)$ with weights $(w_L, w_R)$ using METIS 2-way
\STATE
\STATE $P_L \gets$ RecursiveBisection$(G_L, k_L)$
\STATE $P_R \gets$ RecursiveBisection$(G_R, k_R)$
\STATE
\RETURN $P_L \cup P_R$ (with appropriate label offset for $P_R$)
\end{algorithmic}
\end{algorithm}

\textbf{Key property}: Each binary split is locally optimal (minimizes edge cut between two parts) but doesn't consider final $k$-way partition structure.

\textbf{Tree structure matters}: For $k=7$, splits $[3,4]$, $[4,3]$, $[2,5]$, etc. produce different results despite all reaching 7 partitions. Section~\ref{sec:tree_structure} analyzes this effect.

\subsection{N-Way Partitioning (PartGraphKway)}

Algorithm: Directly partition $G$ into $k$ parts using multilevel k-way refinement.

\begin{algorithm}[H]
\caption{N-Way Partitioning}
\begin{algorithmic}[1]
\STATE \textbf{Input:} Graph $G$, target partitions $k$
\STATE \textbf{Output:} Partition assignment $P: V \to \{1, \ldots, k\}$
\STATE
\STATE Coarsen: $G_0, G_1, \ldots, G_m$ where $|V_m| \approx 100k$
\STATE Initial partition: Create $k$ partitions on $G_m$ (greedy/spectral)
\STATE
\FOR{$i = m-1$ \textbf{down to} $0$}
    \STATE Project partition from $G_{i+1}$ to $G_i$
    \STATE Refine $G_i$ partition using $k$-way Fiduccia-Mattheyses (FM) or greedy refinement
    \STATE \quad \textit{(Considers moves between all $k$ partitions)}
\ENDFOR
\STATE
\RETURN Partition on $G_0$
\end{algorithmic}
\end{algorithm}

\textbf{Key property}: Refinement considers all $k$ partitions simultaneously. A vertex can move from partition $i$ to any partition $j$ if it improves global edge cut and maintains balance.

\subsection{Multi-Constraint Partitioning}

Standard partitioning balances a single vertex weight (population). Multi-constraint partitioning balances multiple weights simultaneously \citep{karypis1999,schloegel2000}.

Given vertex weights $w: V \to \mathbb{R}^c$ (c constraints), each partition $P_i$ must satisfy:
\begin{equation}
|w_j(P_i) - w_{j,target}| \leq \epsilon_j \cdot w_{j,target} \quad \forall j \in \{1, \ldots, c\}
\end{equation}

For redistricting with VRA requirements: $w_v = (p_v, m_v)$ where $p_v$ is population and $m_v$ is minority population.

\textbf{Target weights (tpwgts)}: User specifies desired weight distribution. For $k$ partitions and $c$ constraints, $k \times c$ matrix defines targets. Example: Create 2 MM districts (60\% minority) and 5 non-MM districts from 7 total:
\begin{align*}
\text{Districts 1-2 (MM):} \quad & w_{pop} = 1/7, \quad w_{min} = 0.60/7 \\
\text{Districts 3-7 (non-MM):} \quad & w_{pop} = 1/7, \quad w_{min} = \text{remaining}/5
\end{align*}

Both recursive bisection and n-way support multi-constraint partitioning. Our experiments use 2-constraint optimization to evaluate demographic concentration.

\subsection{Prior Comparisons}

Limited prior work directly compares recursive bisection vs n-way for specific applications:

\textbf{Karypis \& Kumar (1998)} \citep{karypis1998}: Initial METIS paper shows n-way achieves slightly better edge cuts (3-5\%) with comparable runtime. Tested on circuit and finite element graphs, not redistricting.

\textbf{Schloegel et al. (2000)} \citep{schloegel2000}: Multi-constraint partitioning evaluation focuses on scalability, not method comparison. Notes recursive bisection can struggle when constraint distributions are asymmetric.

\textbf{Fiduccia \& Mattheyses (1982)} \citep{fiduccia1982}: Classic FM algorithm for 2-way partitioning. Extended to k-way but greedy nature limits theoretical guarantees.

No prior work systematically evaluates these methods for redistricting where demographic targets create highly asymmetric partition requirements. Our work fills this gap.
